\documentclass[11pt]{amsart}
%prepared in AMSLaTeX, under LaTeX2e
\addtolength{\oddsidemargin}{-1.2in} 
\addtolength{\evensidemargin}{-1.2in}
\addtolength{\topmargin}{-.9in}
\addtolength{\textwidth}{2.0in}
\addtolength{\textheight}{1.5in}

\renewcommand{\baselinestretch}{1.05}

\usepackage{verbatim} % for "comment" environment

\usepackage{palatino}

\usepackage[final]{graphicx}

\usepackage{tikz}
\usetikzlibrary{positioning}

\usepackage{enumitem,xspace,fancyvrb}

\newtheorem*{thm}{Theorem}
\newtheorem*{defn}{Definition}
\newtheorem*{example}{Example}
\newtheorem*{problem}{Problem}
\newtheorem*{remark}{Remark}

\DefineVerbatimEnvironment{mVerb}{Verbatim}{numbersep=2mm,frame=lines,framerule=0.1mm,framesep=2mm,xleftmargin=4mm,fontsize=\footnotesize}

% macros
\usepackage{amssymb}
\newcommand{\bA}{\mathbf{A}}
\newcommand{\bB}{\mathbf{B}}
\newcommand{\bE}{\mathbf{E}}
\newcommand{\bF}{\mathbf{F}}
\newcommand{\bJ}{\mathbf{J}}

\newcommand{\bb}{\mathbf{b}}
\newcommand{\br}{\mathbf{r}}
\newcommand{\bv}{\mathbf{v}}
\newcommand{\bw}{\mathbf{w}}
\newcommand{\bx}{\mathbf{x}}

\newcommand{\Div}{\ensuremath{\nabla\cdot}}
\newcommand{\Curl}{\ensuremath{\nabla\times}}
\newcommand{\eps}{\epsilon}
\newcommand{\grad}{\nabla}
\newcommand{\ip}[2]{\ensuremath{\left<#1,#2\right>}}
\newcommand{\lam}{\lambda}
\newcommand{\lap}{\triangle}

\newcommand{\Null}{\operatorname{null}}
\newcommand{\rank}{\operatorname{rank}}
\newcommand{\range}{\operatorname{range}}
\newcommand{\trace}{\operatorname{tr}}

\newcommand{\RR}{\mathbb{R}}
\newcommand{\ZZ}{\mathbb{Z}}

\newcommand{\prob}[1]{\bigskip\noindent\textbf{#1.}\quad }
\newcommand{\epart}[1]{\bigskip\noindent\textbf{(#1)}\quad }
\newcommand{\ppart}[1]{\,\textbf{(#1)}\quad }

\newcommand{\Matlab}{\textsc{Matlab}\xspace}

\newcommand{\ds}{\displaystyle}

\begin{document}
\begin{center}
  \large
  \sc{Worksheet: Exercises on $H^1$ norms and the Poisson equation}
  \normalsize
\end{center}

\bigskip
\prob{1}  Let $f(x)=\sin(3x)$.

\epart{a} Compute the $H^1(-\pi,\pi)$ norm of $f$ via the original definition.

\vspace{30mm}

\epart{b} Compute the $H^1(S^1)=H_{per}^1(-\pi,\pi)$ norm of $f$ via the Fourier series (summability-of-coefficients) definition.

\vspace{60mm}

\prob{2}  Let $f(x,y)=6y(x-x^2) + 2(y-y^3)$ be the data on the unit square $\Omega=(0,1)^2$.

\epart{a} Show that $u(x,y)=(x-x^2)(y-y^3)$ solves the (strong form) of the homogeneous Dirichlet problem for the Poisson equation on $\Omega$: \quad $\ds -\grad^2 u = f, \quad u=0 \text{ on } \partial \Omega$.

\vspace{70mm}

\epart{b} Describe in a few words, if you can, how I must have come up with this example.

\vfill

\clearpage\newpage
\prob{3}  Let $u(x,y)=\sin(\pi x) e^{\pi y}$ on the unit square $\Omega=(0,1)^2$.

\epart{a} Compute the $H^1$ norm of $u$.

\vspace{100mm}

\epart{b} Show that $u$ solves the (strong form) of the Laplace equation on $\Omega$.

\vspace{70mm}

\epart{c} Describe a boundary value problem, on $\Omega$, which $u$ satisfies.

\vfill

\end{document}
